\documentclass[10pt]{article}
\usepackage{compresr}
\cprSetHeader{Python · LlamaIndex Integration}

\begin{document}

\cprTitleBlock{Compresr SDK · Python}{LlamaIndex Integration}{%
    Compress RAG context and chat memory inside LlamaIndex with a single
    drop-in component. Smaller prompts, lower LLM spend, faster responses,
    no changes to your retrievers or vector store.}

% ===========================================================================
\section{Install}

\begin{shcode}
pip install 'compresr[llamaindex]'
export COMPRESR_API_KEY=cmp_...
\end{shcode}

Requires \cprInlineCode{llama-index-core>=0.11}. The memory block additionally
requires the \cprInlineCode{Memory} API (\cprInlineCode{llama-index-core>=0.12}).

% ===========================================================================
\section{Drop-in node postprocessor}

Add \cprInlineCode{CompresrNodePostprocessor} to
\cprInlineCode{as\_query\_engine()} and every retrieved node is rewritten
with query-aware compression before the response synthesizer sees it.
The query is read automatically from the active query engine.

\begin{pycode}
import os
from llama_index.core import VectorStoreIndex, SimpleDirectoryReader
from compresr.integrations.llamaindex import CompresrNodePostprocessor

documents = SimpleDirectoryReader("docs/").load_data()
index = VectorStoreIndex.from_documents(documents)

compresr_pp = CompresrNodePostprocessor(
    api_key=os.environ["COMPRESR_API_KEY"],
    compression_model="latte_v1",     # query-specific model
    target_compression_ratio=0.5,     # remove ~50% of tokens per node
    min_tokens=200,                   # skip short nodes
)

query_engine = index.as_query_engine(
    similarity_top_k=10,
    node_postprocessors=[compresr_pp],
)

response = query_engine.query(
    "What did the customer say about enterprise pricing?"
)
\end{pycode}

Supports adaptive output budgets via \cprInlineCode{target\_token}.
Fail-open by default: if the Compresr API is unavailable, your pipeline
continues with the original nodes.

% ===========================================================================
\section{Memory block}

Plug \cprInlineCode{CompresrMemoryBlock} into the LlamaIndex 2026
\cprInlineCode{Memory} API to shrink chat history in place — no extra LLM
summarization call, no dropped messages.

\begin{pycode}
import os
from llama_index.core.memory import Memory
from compresr.integrations.llamaindex import CompresrMemoryBlock

memory = Memory.from_defaults(
    session_id="user-42",
    token_limit=8_000,
    memory_blocks=[
        CompresrMemoryBlock(
            api_key=os.environ["COMPRESR_API_KEY"],
            target_token=2_000,     # output budget in tokens
            min_tokens=200,
        )
    ],
)
\end{pycode}

% ===========================================================================
\section{Tool wrapper}

Wrap any \cprInlineCode{FunctionTool} so its return value is compressed
before reaching the agent. Metadata (name, description, schema) is
preserved.

\begin{pycode}
from llama_index.core.tools import FunctionTool
from compresr.integrations.llamaindex import wrap_tool_with_compresr

def search_crm(query: str) -> str:
    """Search the CRM for a customer."""
    return _do_search(query)

tool = FunctionTool.from_defaults(fn=search_crm)
wrapped = wrap_tool_with_compresr(
    tool,
    api_key=os.environ["COMPRESR_API_KEY"],
    target_compression_ratio=0.5,
    query_arg="query",   # use the tool's `query=` arg as compression query
)

from llama_index.core.agent.workflow import FunctionAgent
agent = FunctionAgent(tools=[wrapped], llm=llm)
\end{pycode}

% ===========================================================================
\section{Key parameters}

\renewcommand{\arraystretch}{1.2}

\begin{longtable}{@{}p{4.0cm} p{2.4cm} p{7.7cm}@{}}
\toprule
\textbf{Parameter} & \textbf{Default} & \textbf{Description} \\
\midrule
\endhead

\cprInlineCode{api\_key} & --- &
Your Compresr API key. Or pass a pre-built \cprInlineCode{client}. \\

\cprInlineCode{compression\_model} & \cprInlineCode{"latte\_v1"} &
\cprInlineCode{latte\_v1} (default, query-specific) or
\cprInlineCode{latte\_v2} (faster variant). \\

\cprInlineCode{target\_compression\_ratio} & \cprInlineCode{0.5} &
$[0,1]$ = fraction of tokens to remove (0.5 = remove $\sim$50\%).
$>1$ = Nx factor (4 = $\sim$4x smaller). \\

\cprInlineCode{target\_token} & \cprInlineCode{None} &
Absolute output token budget per node or buffer. When set, overrides
\cprInlineCode{target\_compression\_ratio}. \\

\cprInlineCode{min\_tokens} & \cprInlineCode{200} &
Skip compression when input is below this threshold. Token count is
estimated as $\sim$\,\cprInlineCode{chars / 4} — close enough to filter
small payloads without an extra tokenizer call. \\

\cprInlineCode{coarse} & \cprInlineCode{None} &
\cprInlineCode{True} = paragraph-level (faster).
\cprInlineCode{False} = token-level (finer). \cprInlineCode{None} leaves
the choice to the server's per-model default. \\

\cprInlineCode{query} / \cprInlineCode{query\_arg} /
\cprInlineCode{query\_extractor} & \cprInlineCode{None} &
Override or extract the compression query (tool wrapper). Default lifts
\cprInlineCode{query}, \cprInlineCode{question}, or similar kwargs
automatically. \cprInlineCode{query\_extractor} is a callable
\cprInlineCode{(kwargs: dict) -> str | None} — return \cprInlineCode{None}
to fall back to the auto-default. \\

\cprInlineCode{on\_error} & \cprInlineCode{"passthrough"} &
\cprInlineCode{"passthrough"} = log and return original content on
failure. \cprInlineCode{"raise"} = propagate. \\
\bottomrule
\end{longtable}

\end{document}
